\cotefile{136}
\scanpage{64}

% top, opening two lines
Mais si $F$ \inserted{sur $X_*$} est \deleted{loc. constant}\inserted{un $k$-Module} \inserted{($k=$ un anneau \unsure{constant comm.|commutatif})}, alors on vérifie facilement que

% upper, first display and its underbrace
\[
\unsure{C^{*}|C^{\bullet}}(X_*,F)\simeq\Gamma\bigl(X_*,\underbrace{(C^{*}_{*}\mid\unsure{X_{*}|XX_{*}})}_{C^{*}_{X_*}}\otimes_{k_*}F\bigr)
\]

% left, encircled beside first display
\margin{NB On pose $C^{*}_{*}=C^{*}_{*k}$ pour abréger}

% upper middle, resolution, marginal query and conclusion
et on vérifie tout de suite que $C^{*}_{X_\bullet}\otimes_{k_\bullet}F$ est une résolution de $F$, laquelle est d'ailleurs \uline{flasque}\margin{à vérifier ce point $\longrightarrow$} \deleted{Donc} (que $F$ soit loc. constant ou non, d'ailleurs). Donc, sans hypothèse de locale constance sur $F$, $C^{*}(X_\bullet,F)$ permet de calculer la coh. de $X_\bullet$ à coéff. dans $F$ (i.e. on a un isom. can.

% middle, cohomology display and definition beneath C
\[
H^{*}\bigl(\underbrace{C^{*}(X_*,F]}_{\substack{\Vert\,\mathrm{dfn}\\\Gamma(X_\bullet,C^{*}_{X_\bullet}\otimes_k F}}\bigr)\simeq H^{*}(X_\bullet,F)
\]

% middle to lower, functoriality and hypercohomology
pour \deleted{\unsure{compatible}} fonctoriel en $F$ (et $=$ en $(X_\bullet,F)\ldots$) et compatible avec les hom. cobords. De $=$ pour le calcul de l'hypercohomologie de $X_\bullet$ à coeff. dans un complexe de $k$-Modules $F^{*}$. En termes de catégories dérivées, on aura un isomorphisme de foncteurs triangulés

% bottom, derived functor and explanation below brace
\[
R\Gamma(X_*,F^*)\simeq\underbrace{C^{*}(X_*,F^*)}_{\substack{\text{pour degré total}\\\text{associé aux}\\\text{deux degrés}\\\text{(\unsure{prov.|proven.|venant} de \unsure{\v{C} et|\v{C}ech|C\textsuperscript{*} et} de}\\ F^{*}\text{)}}}
\]

% lower left
% archive sheet number: 63
